\documentclass[12pt,a4paper]{article}
\usepackage[utf8]{inputenc}
\usepackage[T1]{fontenc}
\usepackage[ngerman]{babel}
\usepackage{amsmath,amssymb,amsthm}
\usepackage{geometry}
\usepackage{titlesec}
\usepackage{fancyhdr}
\usepackage{mathrsfs}

\geometry{margin=2.5cm}

\pagestyle{fancy}
\fancyhf{}
\fancyhead[C]{\small Lehrbuch der Algebra --- Band 3}
\fancyfoot[C]{\thepage}
\renewcommand{\headrulewidth}{0.4pt}

\title{Lehrbuch der Algebra --- Band 3, Teil 13}
\author{Heinrich Weber}
\date{}

\newtheorem*{satz}{Satz}
\newtheorem*{bem}{Bemerkung}

\begin{document}
\maketitle
\thispagestyle{fancy}

\section*{\S 152. Die Pole der Funktion $\wp(uw)$.}
\addcontentsline{toc}{section}{\S 152. Die Pole der Funktion $\wp(uw)$.}

worin $v$ ebenso wie $u$ eine ganze Zahl des K"orpers $\Omega$ ist, die nicht durch $u$ teilbar ist.

Ist diese Bedingung erf"ullt, so ist $\varrho$ eine Wurzel von $P_u(x)$.

Weil aber $\wp(vu)$ in der $u$-Ebene f"ur
\[v \equiv 0 \pmod{\omega_1, \omega_2}\]
unendlich in der zweiten Ordnung wird, so ist $P_u(x)$ durch $(x - \varrho)^2$ teilbar, es sei denn, da\ss{} $\wp(v) - \varrho$ f"ur $v = v{:}u$ selbst in der zweiten Ordnung verschwindet; dies tritt wegen \S 46, (18) dann ein, wenn
\[\frac{v}{u} = \frac{\omega}{2}\]
ist, also wenn $v{:}u$ eine halbe Periode ist, und also
\begin{equation}\varrho = e_1, e_2, e_3\end{equation}
wird.

1. Nach (4) ist, wenn $\omega_1 = 1$, $\omega_2 = \omega$ gesetzt ist, jede ganze Zahl des K"orpers $\Omega$ eine Periode von $\wp(u)$. Dies gilt nicht umgekehrt; da aber $A\omega$ eine ganze Zahl in $\Omega$ ist, so mu\ss{} jede Periode durch Multiplikation mit $A$ in eine ganze Zahl verwandelt werden. Daraus folgt, da\ss{} zwei Werte
\[\wp\left(\frac{v}{u}\right), \quad \wp\left(\frac{v'}{u}\right)\]
dann und nur dann einander gleich sind, wenn
\[v \equiv \pm v' \pmod{u}\]
ist.

Denn ist $\varrho = \varrho'$, so mu\ss{} entweder
\[\frac{v-v'}{u} \quad \text{oder} \quad \frac{v+v'}{u}\]
durch Multiplikation mit $A$ in eine ganze Zahl verwandelt werden. $A$ ist aber relativ prim zu $u$, und folglich mu\ss{} eine dieser beiden Zahlen selbst ganz sein.

Die Anzahl der inkongruenten Werte von $v$ ist aber gleich $N(u) = m$ (Bd.~II, \S 165). Lassen wir also $v$ ein volles Restsystem nach dem Modul $u$ mit Ausschlu\ss{} der Null durchlaufen, so bekommen wir jeden Wert von $\varrho$ zweimal, au\ss{}er wenn
\begin{equation}2v \equiv 0 \pmod{u},\end{equation}
und in diesem Falle ist $\wp\left(\dfrac{v}{u}\right)$ einer der Werte $e$. Daraus folgt:
\begin{equation}P_u(x) = C\prod_{v} \left\{x - \wp\left(\frac{v}{u}\right)\right\}^2,\end{equation}
und $P_u(x)$ ist vom Grade $m-1$.

In dem Produkt $P_u(x)$ kommt jeder Faktor $x - \varrho$ zweimal vor, au\ss{}er wenn $\varrho$ einem der $e$ gleich ist, wenn also die Kongruenz (7) f"ur ein von Null verschiedenes $v$ erf"ullt werden kann.

1) Wenn $u$ relativ prim zu 2 ist, also wenn $m$ ungerade ist, so hat die Kongruenz (7) nur die Wurzel $v = 0$. Dann ist $P_u(x)$ ein Quadrat, und wenn $\Phi(x)$ eine Funktion vom Grade $\frac{1}{2}(m-1)$ ist,
\begin{equation}P_u(x) = \Phi(x)^2.\end{equation}

Wenn $u$ mit 2 einen Teiler gemein hat, sind drei F"alle zu unterscheiden.

2) Wenn $u$ durch 2 teilbar ist, so mu\ss{} $v$, wenn es der Kongruenz (7) gen"ugen soll, ein Vielfaches von $\frac{u}{2}$ sein, und da es drei von Null verschiedene nach dem Modul 2 inkongruente Zahlen in $\Omega$ gibt, so hat (7) drei Wurzeln. Es kommen also unter den $\varrho$ die drei Werte $e_1, e_2, e_3$ vor, und wir haben:
\begin{equation}P_u(x) = \Phi(x)^2 \cdot \Psi, \qquad m \equiv 0 \pmod{2},\end{equation}
worin $\Psi$ eine ganze Funktion vom Grade $\frac{1}{2}m - 2$ ist.

3) Wenn 2 im K"orper $\Omega$ in zwei (gleiche oder verschiedene) Primideale zerf"allt, also wenn
\[\Delta \equiv 0 \pmod{4} \quad \text{oder} \quad \equiv 1 \pmod{8}\]
ist, so kann $u$ durch einen dieser Primfaktoren teilbar sein, ohne durch 2 teilbar zu sein. Dann hat die Kongruenz (7) nur eine von Null verschiedene Wurzel und es kommt unter den Zahlen $\varrho$ nur eine der drei Zahlen (6), etwa $e_1$, vor. In diesem Falle ist
\begin{equation}P_u(x) = [x - \wp(u) - e_1] \cdot \Phi^2,\end{equation}
worin $\Phi$ eine Funktion vom Grade $\frac{1}{2}(m-1)$ ist.

In bezug auf die beiden besonderen F"alle $g_2 = 0$ und $g_3 = 0$ ist noch folgendes zu bemerken.

4) Wenn $g_2 = 0$ ist, so ist $\Delta = -4$, $j(\omega) = 0$, und es ergibt sich aus der Differentialgleichung \S 46, (18):
\[\wp'^2(u) = 4\wp^3(u) - g_3\wp(u),\]
durch die mit R"ucksicht auf das Anfangsglied der Entwickelung die $\wp$-Funktion eindeutig bestimmt ist.

Ersetzt man dann $u$ durch $\varepsilon u$, so folgt:
\[-\wp'(u)^2 = 4\wp^3(u) - g_3\wp(u),\]
und daraus:
\[\wp(\varepsilon u) = \varepsilon\wp(u).\]

Die Gr"o\ss{}en $\varrho$ sind also einander paarweise entgegengesetzt, eine der drei Gr"o\ss{}en $e_1, e_2, e_3$ ist gleich Null, die beiden anderen ebenfalls entgegengesetzt gleich. Daraus folgt, da\ss{} in den Formeln (9) bis (11) die Funktion $\Phi$ in diesem Falle nur die geraden Potenzen von $\wp(u)$ enthalten kann.

5) Ist $g_3 = 0$, so ist $\Delta = -3$, $j(\omega) = 64 \cdot 27$, und es ist, wenn mit $\varrho$ eine imagin"are dritte Einheitswurzel bezeichnet wird:
\begin{equation}\wp(\varrho u) = \varrho\wp(u) = \varrho'\wp(\varrho^2 u),\end{equation}
und die Wurzeln $\varrho$ von $\Phi$ ordnen sich zu dreien in der Weise:
\begin{equation}\varrho, \quad \varrho\varrho', \quad \varrho\varrho'^2.\end{equation}

Unter diesen drei Werten sind nur dann zwei einander gleich, wenn sie alle drei verschwinden. Aus (13) aber ergibt sich, da\ss{} dies eintritt, wenn $v$ mit einem der Werte $\pm 1 \pm \sqrt{-3}$ kongruent wird. Denn f"ur $\Delta = -3$ sind $1, \varrho, \varrho^2$ Perioden von $\wp(u)$, und da nun
\[\frac{v}{u}\sqrt{-3}\]
ist, so ist nach (15):
\[\wp\left(\frac{v}{u}\right) = 0,\]
und der Wert Null kommt also unter den $\varrho$ dann und nur dann vor, wenn $m$ durch 3 teilbar ist. Daraus folgt:

Ist $\Delta = -3$ und $m \equiv 0 \pmod{3}$, so ist $\Phi\cdot\wp(u)$ eine rationale Funktion von $\wp(u)$.

Ist $\Delta = -3$ und $m \equiv 1 \pmod{3}$, so ist $\Phi$ selbst eine rationale Funktion von $\wp(u)^3$.

[Da $m$ die Form $4(x^2 + 3y^2)$ haben mu\ss{}, so kann hier $m$ nicht $\equiv -1 \pmod{3}$ sein.]

\section*{\S 153. Die Funktion $z(u)$.}
\addcontentsline{toc}{section}{\S 153. Die Funktion $z(u)$.}

Wir f"uhren nun nach \S 150 (6), (8), (9), (10), (12), (13) eine Funktion $z = r(u)$ ein, die nur von zwei Argumenten $\omega$, $\omega'$ abh"angt, indem wir setzen:

a) im allgemeinen,
\begin{equation}z = \wp'(u)^2,\end{equation}
b) wenn $g_2 = 0$,
\[z = \wp(u)^2,\]
c) wenn $g_3 = 0$,
\[z = \wp(u)^3.\]

Und wenn dann $u$ wie oben eine ganze Zahl des K"orpers $\Omega$ ist, so k"onnen wir in allen drei F"allen setzen:
\begin{equation}t(u) = \frac{Z(z)}{N(z)},\end{equation}
worin $Z$, $N$ ganze Funktionen von $z$ sind, die bis auf konstante Faktoren im Falle (1) a) mit $P$, $\Phi$ des vorigen Paragraphen "ubereinstimmen, in den F"allen (1) b) und (1) c) daraus durch Erhebung ins Quadrat oder in den Kubus hervorgehen [\S 152, 4), 5)].

Wir richten den Bruch (2) so ein, da\ss{} die h"ochste Potenz von $z$ im Nenner $N$ den Koeffizienten 1 hat. Die "ubrigen Koeffizienten von $N$ sind dann rational zusammengesetzt aus Gr"o\ss{}en der Form
\[\varrho = \wp\left(\frac{v}{u}\right)\]
und da $v/u$ eine Periode von $t$ ist, so geh"oren diese Gr"o\ss{}en zu den Wurzeln der Teilungsgleichung der Perioden, und sind daher, da der Modul der entsprechenden elliptischen Funktionen eine algebraische Zahl ist, selbst algebraische Zahlen (\S 58, \S 61). Daraus folgt:

\begin{satz}
Die Koeffizienten in $N(z)$ sind algebraische Zahlen.
\end{satz}

Wir erweitern den Bruch $Z:N$ durch Multiplikation von Z"ahler und Nenner mit dem Produkt der konjugierten Werte $N'$, $N''$, $N'''$, \ldots und erhalten:
\begin{equation}t(u) = \frac{Z}{N},\end{equation}
worin nun $Z$ und $N$ ganze Funktionen von $z$ sind, die einen gemeinsamen Teiler enthalten k"onnen, wobei jedoch $N$ rationale Zahlenkoeffizienten hat.

\section*{\S 154. Der Teilungsk\"orper.}
\addcontentsline{toc}{section}{\S 154. Der Teilungsk\"orper.}

Nach \S 150, (8), (10), (13) beginnt die Entwickelung von $t(u)$ nach steigenden Potenzen von $u$ mit einer negativen Potenz von $u$.

Wir ordnen in der Gleichung
\begin{equation}Z(u) - t(u) \cdot N = 0\end{equation}
die Funktion $Z$ nach fallenden Potenzen von $z$, entwickeln beide Seiten nach steigenden Potenzen von $u$ und vergleichen die Koeffizienten gleich hoher Potenzen von $u$. Dann bekommen wir f\"ur die Bestimmung der Koeffizienten von $Z$ eine Reihe linearer Gleichungen, deren jede folgende nur einen neuen dieser Koeffizienten enth\"alt, und daraus k\"onnen diese Koeffizienten successiv rational berechnet werden.

Beachtet man nun die Form der Entwickelungen des \S 150, so ergibt sich, da\ss{} die Koeffizienten von $Z$ rationale Funktionen von $j(\omega)$ und $\sqrt{A}$ sind. Befreit man nach dem Euklidischen Algorithmus $N$ und $Z$ von gemeinschaftlichen Faktoren, so kommt man zu den Funktionen $P$, $\Phi$ zur\"uck, und es ergibt sich, wenn wir unter dem Klassenk\"orper den Inbegriff der rationalen Funktionen von $\sqrt{A}$ und $j(\omega)$ verstehen:

\begin{satz}
Die Koeffizienten der Funktionen $\Phi(z)$, $P(z)$ in (2) geh\"oren dem Klassenk\"orper an.
\end{satz}

Wenn wir (durch rationale Rechnung) die Funktion $\Phi(x)$ von allen mehrfach darin vorkommenden Faktoren befreien, so bleibt eine ganze Funktion $\Psi_u(x)$ von $x$ u\"brig, deren Koeffizienten dem Klassenk\"orper angeh\"oren, und die Wurzeln von $\Psi_u(x)$ sind s\"amtliche voneinander verschiedene unter den Gr\"o\ss{}en
\[r = t\left(\frac{v}{u}\right).\]

Zwei dieser Gr\"o\ss{}en $r$, $r'$ sind einander gleich, wenn im Falle (1) a), wenn
\[\frac{v}{u} = \pm\frac{v'}{u},\]
im Falle (1) b), wenn
\[\frac{v}{u} = \pm\frac{v'}{u}, \quad \pm i\frac{v'}{u},\]
im Falle (1) c), wenn
\[\frac{v}{u} = \pm\frac{v'}{u}, \quad \pm \varrho\frac{v'}{u}, \quad \pm \varrho^2\frac{v'}{u},\]
worin die Kongruenz auf die Perioden bezogen wird, also (\S 152, 1.) im Falle a), wenn
\[v \equiv \pm v' \pmod{u},\]
b), wenn
\[v \equiv \pm v', \quad \pm iv' \pmod{u},\]
c), wenn
\[v \equiv \pm v', \quad \pm \varrho v', \quad \pm \varrho^2 v' \pmod{u}.\]

Es seien jetzt
\begin{equation}\mu, \quad \mu_1\end{equation}
zwei ganze Zahlen mit dem gr\"o\ss{}ten gemeinschaftlichen Idealteiler $m$ im K\"orper $\Omega$.

Dann haben $\Psi_u$ und $\Psi_{u_1}$ einen gemeinsamen Linearteiler, wenn
\[r = r',\]
und dies findet nach (2) dann und nur dann statt, wenn
\begin{equation}uv' \equiv \pm \varepsilon u'v \pmod{um},\end{equation}
worin $\varepsilon$ eine Einheit des K\"orpers $\Omega$ ist, also $\varepsilon = \pm 1$ im allgemeinen, $\varepsilon = \pm 1$, $\pm i$, wenn $\Delta = -4$ und $\varepsilon = \pm 1$, $\pm \varrho$, $\pm \varrho^2$ wenn $\Delta = -3$ ist.

Aus (5) folgt aber, da\ss{} $v$ durch $a$ und $v_1$ durch $a_1$ teilbar sein mu\ss{}.

Denken wir uns $v$ gegeben und $v_1$ gesucht, so ist die Kongruenz (5) nur m\"oglich, wenn $u_1v$ durch $u$ teilbar ist, und da $a$ und $a_1$ relativ prim sind, so mu\ss{} $v$ durch $a$ teilbar sein.

Ist $\alpha$ eine durch $a$ teilbare ganze Zahl in $\Omega$, so beschaffen, da\ss{} $\alpha:a$ relativ prim zu $m$ ist, so k\"onnen wir demnach
\begin{equation}v = \alpha \xi \pmod{u}\end{equation}
setzen (Bd.~II, \S 166, 7.) und erhalten die s\"amtlichen voneinander verschiedenen Werte $t\left(\frac{v}{u}\right)$ und jeden zwei-, vier- oder sechsmal, wenn wir $\xi$ ein volles Restsystem nach dem Modul $m$ durchlaufen lassen (mit Ausschlu\ss{} der Null). Man erh\"alt dann die s\"amtlichen gemeinsamen Wurzeln von $\Psi_u$ und $\Psi_{u_1}$ in der Form
\begin{equation}r = t\left(\frac{\alpha \xi}{u}\right).\end{equation}

Der gr\"o\ss{}te gemeinschaftliche Teiler $D_u$ von $\Psi_u$ und $\Psi_{u_1}$ hat also alle die Gr\"o\ss{}en (7) zu Wurzeln.

Aus $D_u$ l\"a\ss{}t sich noch auf rationalem Wege ein Teiler $\Psi_m$ absondern, der nur die unter den Gr\"o\ss{}en (7) zu Wurzeln hat, in denen $\xi$ relativ prim zu $m$ ist, wenn man $D_u$ von allen Faktoren befreit, die es mit einem $D_{m'}$ gemein hat, wenn $m'$ ein Teiler von $m$ ist.

Da in dieser Betrachtung in (3) $m$ jedes beliebige Ideal in $\Omega$ bedeuten kann, so kommen wir zu folgendem Hauptsatz:

\begin{satz}
Ist $m$ ein beliebiges Ideal des K\"orpers $\Omega$, so existiert eine Funktion $\Psi_m$ in $z$, deren Wurzeln die voneinander verschiedenen der Zahlen $r_\xi$ sind, wenn $\xi$ ein System inkongruenter, zu $m$ teilerfremder Zahlen durchl\"auft.
\end{satz}

Um den Grad der Funktion $\Psi_m$ zu bestimmen, bemerken wir, da\ss{} zwei Gr\"o\ss{}en $r_\xi$, $r_{\xi'}$ nur dann einander gleich sind, wenn
\[\xi \equiv \pm \varepsilon \xi' \pmod{m}\]
ist.

Es werden also so viele von den $r_\xi$ einander gleich, als es modulo $m$ inkongruente Einheiten gibt; das sind im allgemeinen zwei, f\"ur $\Delta = -4$ sind es vier und f\"ur $\Delta = -3$ sind es sechs. Diese Zahlen verringern sich aber wiederum, wenn verschiedene der Einheiten $\varepsilon$ nach dem Modul $m$ kongruent sind, wenn also $1 - \varepsilon$ durch $m$ teilbar ist. Das kann aber nur vorkommen, wenn $m$ ein Teiler von 2 oder von 3 ist. Also:

\begin{satz}
Der Grad der Funktion $\Psi_m$ ist gleich der Anzahl $\psi(m)$ (Bd.~II, \S 168) der nach dem Modul $m$ inkongruenten, zu $m$ teilerfremden Zahlen in $\Omega$, geteilt durch die Anzahl der nach $m$ inkongruenten Einheiten in $\Omega$.
\end{satz}

Die Wurzeln der Funktion $\Psi_m$ bestimmen einen algebraischen K\"orper $\mathfrak{T}_m$ \"uber dem Klassenk\"orper, dessen relativer Grad h\"ochstens gleich dem Grade von $\Psi_m$ ist, und beide Grade sind gleich, wenn $\Psi_m$ irreduzibel ist.

\begin{satz}
Diesen K\"orper $\mathfrak{T}_m$ wollen wir den Teilungsk\"orper f\"ur den Modul $m$ nennen.
\end{satz}

Diese Teilungsk\"orper spielen f\"ur den K\"orper $\Omega$ eine \"ahnliche Rolle, wie die Kreisteilungsk\"orper f\"ur den K\"orper der rationalen Zahlen, nur da\ss{} sich hier noch der Klassenk\"orper dazwischen schiebt, zu dem es im K\"orper der rationalen Zahlen, ebenso wie in jedem einklassigen K\"orper, kein Analogon gibt.

Nach dem Multiplikationstheorem (\S 151) kann
\begin{equation}t(\xi u) = R_\xi(t(u))\end{equation}
rational durch $t(u)$ ausgedr\"uckt werden.

Setzen wir
\[u = \frac{\xi}{m} = \frac{\lambda}{k},\]
so folgt:
\begin{equation}t(\lambda) = R_\xi(t(\xi)),\end{equation}
und folglich durch Vertauschung von $\xi$ mit $\xi'$:
\begin{equation}t(\xi') = R_{\xi'}(t(\lambda)),\end{equation}
und damit nach Bd.~I, \S 169:

\begin{satz}
Der Teilungsk\"orper ist in bezug auf den Klassenk\"orper relativ Abelsch.
\end{satz}

Die Zahlen $\xi$, nach dem Modul $m$ genommen, bilden bei der Multiplikation eine Gruppe, und die Relativgruppe des Teilungsk\"orpers ist mit dieser Gruppe isomorph (oder wenigstens mit einem Teiler dieser Gruppe. Die Irreduzibilit\"at von $\Psi_m$ wird weiterhin bewiesen werden, wodurch dann die Gruppe von $\mathfrak{T}_m$ selbst festgestellt ist).

\section*{\S 155. Multiplikation der elliptischen Funktionen f\"ur einen ungeraden Multiplikator.}
\addcontentsline{toc}{section}{\S 155. Multiplikation der elliptischen Funktionen f\"ur einen ungeraden Multiplikator.}

F\"ur den Nachweis der Existenz des Teilungsk\"orpers ist die Benutzung der Weierstrassschen Funktion und der daraus abgeleiteten $t$-Funktion sehr zweckm\"a\ss{}ig, und es w\"are am bequemsten, wenn man darauf auch die weitere Untersuchung dieses K\"orpers gr\"unden k\"onnte. Daf\"ur aber ist die arithmetische Natur der Multiplikationsformeln noch nicht gen\"ugend bekannt, und es ist darum zurzeit noch notwendig, die komplexe Multiplikation und Teilung auch der Jacobischen elliptischen Funktionen zu betrachten.

Hierbei wollen wir den Multiplikationsformeln, die wir in \S 57 abgeleitet haben, noch eine etwas andere Gestalt geben.

Wir betrachten zun\"achst durchweg den Multiplikator $m$ als ungerade Zahl.

Wir setzen, wenn
\begin{equation}x = \sqrt{k} = \frac{\vartheta_2(\omega)}{\vartheta_3(\omega)}\end{equation}
den Modul der elliptischen Funktionen bedeutet, nach Kronecker\footnote{Zur Theorie der elliptischen Funktionen. Sitzungsbericht der Berliner Akademie vom 29. Juli 1886.}:
\begin{equation}A = A(x) = x^2,\end{equation}
und entwickeln diesen Ausdruck nach steigenden Potenzen von
\[q = e^{\pi i \omega}.\]

Man kann die Form dieser Entwickelung aus den Produktformeln [\S 24, (11)] ableiten, und findet:
\begin{equation}A = A_0 + A_1 q + A_2 q^2 + A_3 q^3 + \cdots,\end{equation}
worin die Konstanten $A_0, A_1, A_2, \ldots$ ganze rationale Zahlen sind.

Setzen wir ferner
\begin{equation}x = \sqrt{k} \, \operatorname{sn}(2Kv, k^2),\end{equation}
so bleibt (4) nach \S 45, (9) unge\"andert, wenn $x$ mit $1:x$ vertauscht wird.

Es gen\"ugt $x$ der Differentialgleichung
\begin{equation}\left(\frac{dx}{dv}\right)^2 = 1 - (1 + A)x^2 + Ax^4.\end{equation}

Entwickelt man also $x$ nach dem Taylorschen Lehrsatz in eine Reihe nach steigenden Potenzen von $v$:
\begin{equation}x = v + X_1 v^3 + X_2 v^5 + \cdots,\end{equation}
so sind die Koeffizienten $X_1, X_2, \ldots$ ganze rationale Funktionen von $A$ mit rationalen Zahlenkoeffizienten.

Benutzen wir die Bezeichnung von \S 42, so ist
\begin{equation}x = \frac{\sigma_1}{\sigma_0} \cdot \frac{\sigma(v)}{\sigma_3(v)} \cdot e^{-(e_1-e_3)uv},\end{equation}
und nach \S 57, IV. mit einer etwas modifizierten Bedeutung der $A, B, C, D$:
\begin{equation}y_1 \operatorname{sn}(mv) = \frac{x A(x)}{D(x)}, \quad y_1 \operatorname{cn}(mv) = \frac{B(x)}{D(x)}, \quad y_1 \operatorname{dn}(mv) = \frac{C(x)}{D(x)}.\end{equation}

Hier sind $A(x)$, $B(x)$, $C(x)$, $D(x)$ ganze rationale Funktionen von $x$ vom Grade $\frac{1}{2}(m^2-1)$, und es ist aus den Rekursionsformeln \S 57, (13), (14) zu ersehen, da\ss{} die Funktionen $A$, $D$ und das Produkt $BC$ ganze rationale Funktionen von $A$ sind, deren Zahlenkoeffizienten rational sind.

Da\ss{} es ganze Zahlen sind, ist wegen des Nenners 2, der zun\"achst auftritt, nicht zu ersehen.

Um dies n\"aher zu untersuchen, betrachten wir die Wurzeln der Funktion $A(x)$:

Sie sind, wenn gesetzt wird:
\[z_{h,h'} = \frac{2hK + 2h'iK'}{m}, \quad h, h' = 0, 1, \ldots, m-1,\]
worin $h,h'$ irgend welche ganze Zahlen sein k\"onnen; nur d\"urfen sie nicht beide gleich Null sein.

Man erh\"alt alle Wurzeln von $A(x)$, wenn man $h, h'$ je ein volles Restsystem nach dem Modul $m$ durchlaufen l\"a\ss{}t, abgesehen von der Kombination $0,0$.

Es ist dann
\begin{equation}x_{h,h'} = \sqrt{k} \, \operatorname{sn} \frac{2hK + 2h'iK'}{m}.\end{equation}

Um diesen Ausdruck nach steigenden Potenzen von $q$ zu entwickeln, mu\ss{} man zun\"achst die niedrigste Potenz von $q$ im Nenner aufsuchen, d.~h. das Glied, in dem $v$ so klein als m\"oglich wird. Dies findet statt, wenn $v$ die zun\"achst an $-h'/m$ gelegene ganze Zahl ist. Es gibt, da $m$ ungerade ist, nur eine solche Zahl $v$, und die Formel (10) erh\"alt die Gestalt:
\begin{equation}x_{h,h'} = \frac{P_1}{Q_1},\end{equation}
worin $Q_1$ und $P_1$ nach steigenden Potenzen von $q$ geordnete Reihen sind, deren Koeffizienten ganze algebraische Zahlen (Kreisteilungszahlen) sind.

Demnach l\"a\ss{}t sich $x_{h,h'}$ nach steigenden Potenzen von $q$ entwickeln und die Koeffizienten dieser Entwickelung sind ganze algebraische Zahlen.

Ist $S$ eine symmetrische Funktion der $x_{h,h'}$ mit rationalen ganzzahligen Koeffizienten, so l\"a\ss{}t auch diese sich in derselben Weise nach Potenzen von $q$ entwickeln. Andererseits ist $S$ nach dem Fundamentalsatz \"uber symmetrische Funktionen rational durch die Koeffizienten von $A(x)$ ausdr\"uckbar, ist also eine ganze rationale Funktion von $A$ mit rationalen Zahlenkoeffizienten von der Form:
\begin{equation}S = S_0 + S_1 A + S_2 A^2 + S_3 A^3 + \cdots,\end{equation}
worin die $S_0, S_1, S_2, \ldots$ rationale Zahlen sind. Da\ss{} es ganze Zahlen sind, soll eben bewiesen werden.

Das ergibt sich aber sehr einfach, wenn man in (12) f\"ur $A$ die Entwickelung (3) einsetzt und dann die Koeffizienten gleich hoher Potenzen von $q$ auf beiden Seiten miteinander vergleicht. Man bekommt n\"amlich zur Bestimmung der $S_0, S_1, S_2, \ldots$ eine Reihe linearer Gleichungen, deren jede folgende nur ein neues $S_i$ und zwar mit dem Koeffizienten 1 enth\"alt, w\"ahrend auf der anderen Seite dieser Gleichungen ganze algebraische Zahlen stehen. Hiernach sind also die $S_0, S_1, S_2, \ldots$ ganze algebraische Zahlen, und da sie rational sind, sind es auch ganze rationale Zahlen.

Da wir \"uberdies nach \S 57, II und (7) $A(0)$ und $A(\infty)$ kennen, so ergibt sich f\"ur $A(x)$ der Ausdruck:
\begin{equation}A(x) = \pm x \prod (x^2 - x_{h,h'}^2),\end{equation}
worin die $x_{h,h'}^2$ ganze ganzzahlige Funktionen von $A$ sind.

Ferner ist nach \S 57, (7):
\[A(x) D(x) = x^{m^2} + \cdots,\]
und folglich:
\begin{equation}\pm D(x) = \pm 1 + A_1 x^2 + A_2 x^4 + \cdots + m^2 x^{m^2-1}.\end{equation}

[Das Zeichen $\pm$ in (13) und (14) ist nach \S 57 bestimmt durch $m \equiv \pm 1 \pmod{4}$.]

F\"ur die Funktionen $B$ und $C$ gilt die Relation:
\begin{equation}B(x) C(x) = x^{m^2-1} + \cdots,\end{equation}
und ihre Wurzeln sind daher zueinander reziprok. Der erste und der letzte Koeffizient sind in $B$ und in $C$ gleich der Einheit.

\section*{\S 156. Ubergang zu den singul\"aren Moduln.}
\addcontentsline{toc}{section}{\S 156. Ubergang zu den singul\"aren Moduln.}

Wir nehmen jetzt an, da\ss{} $\omega$ die Wurzel einer quadratischen Gleichung:
\begin{equation}a\omega^2 + b\omega + c = 0\end{equation}
mit der negativen Stammdiskriminante
\begin{equation}\Delta = b^2 - 4ac\end{equation}
sei.

Ist dann $j(\omega)$ die Invariante, so gen\"ugt die Gr\"o\ss{}e $A$ [\S 155, (2)] der Gleichung 6ten Grades:
\begin{equation}A^6 - u_1 A^5 + u_2 A^4 - [j(\omega) - 27 \cdot 64] A^3 + 64 [j(\omega) - 27 \cdot 64] A^2 - 64^2 u_1 A + 64^3 = 0,\end{equation}
deren Koeffizienten ganze algebraische Zahlen sind. Folglich ist $A$ selbst eine ganze algebraische Zahl (Bd.~II, \S 154, 11.)

und in \S 135 haben wir gesehen, da\ss{} $A$ Klasseninvariante der Diskriminante $4\Delta$ oder $16\Delta$ ist.

Wir adjungieren also dem Klassenk\"orper
\begin{equation}\Omega = \Omega(j(\omega))\end{equation}
die Zahl $A$ und erhalten einen Klassenk\"orper:
\begin{equation}\mathfrak{K} = \Omega(A^2), \quad \text{wenn } \Delta \equiv 0 \pmod{4}, \text{ oder } \equiv 1 \pmod{8},\end{equation}
\[\mathfrak{K} = \Omega(A), \quad \text{wenn } \Delta \equiv 5 \pmod{8}.\]

Ist $h$ die Klassenzahl von $\mathfrak{f}$, also der Relativgrad von $\Omega(\mathfrak{f})$ in bezug auf den quadratischen K\"orper $\Omega = \mathbf{R}(\sqrt{\Delta})$, und $W$ der Relativgrad von $\mathfrak{K}$ in bezug auf $\Omega$, so ist (\S 100, \S 123):
\begin{equation}W = h, \text{ wenn } \Delta \equiv 1 \pmod{8},\end{equation}
\[W = 2h, \text{ wenn } \Delta \equiv 0 \pmod{4}, \text{ oder } \equiv 5 \pmod{8}.\]

In dem ersten dieser drei F\"alle ist also $\mathfrak{K}$ mit $\Omega(\mathfrak{f})$ identisch.

Durch die Substitution des singul\"aren Moduls $\omega$ gehen die Koeffizienten $a_1, a_2, \ldots$; $d_1, d_2, \ldots$ in \S 155, (13), (18) in ganze algebraische Zahlen \"uber, und daraus folgt, Bd.~II, \S 154, 11., da\ss{} die Wurzeln von $A(x)$, n\"amlich:
\begin{equation}\sqrt{k} \, \operatorname{sn} \frac{2hK + 2h'iK'}{m}\end{equation}
ganze Zahlen und die Wurzeln von $B(x)$:
\begin{equation}\sqrt{k} \, \operatorname{cn} \frac{2hK + 2h'iK'}{m}\end{equation}
Einheiten sind.

\section*{\S 157. Komplexe Multiplikatoren.}
\addcontentsline{toc}{section}{\S 157. Komplexe Multiplikatoren.}

Es gen\"uge $\omega$ wie im vorigen Paragraphen der quadratischen Gleichung
\begin{equation}a\omega^2 + b\omega + c = 0\end{equation}
mit negativer Stammdiskriminante
\begin{equation}\Delta = b^2 - 4ac,\end{equation}
\[\omega = \frac{-b + \sqrt{\Delta}}{2a}.\]

Wir legen unseren Betrachtungen wie im vorigen Paragraphen die Funktion
\begin{equation}y_1 \operatorname{snv} = \sqrt{k} \, \frac{\sigma(v)}{\sigma_3(v)} = \frac{s(u)}{S(u)}\end{equation}
zugrunde, die wir als Funktion von $u = v/2K$ mit $S(u)$ bezeichnen.

Diese Funktion hat dann die durch
\begin{equation}S(u+1) = -S(u), \quad S(u+\omega) = -S(u)\end{equation}
ausgedr\"uckte doppelte Periodizit\"at.

Alle Perioden von $S(u)$ sind in der Form $2m + n\omega$ enthalten, worin $m, n$ ganze rationale Zahlen sind, und gehen also durch Multiplikation mit $\omega$ in ganze Zahlen \"uber.

Zwei Zahlen $u, w$, die sich nur um eine Periode unterscheiden, hei\ss{}en kongruent nach dem Modul $2, \omega$:
\[w \equiv u \pmod{2, \omega}.\]

Da die Funktion $\operatorname{sn} v$ einen Wert nur zweimal im Periodenparallelogramm annimmt, so wird nur dann $S(u) = S(w)$, wenn entweder
\begin{equation}w \equiv u \pmod{2, \omega}\end{equation}
oder
\[w \equiv 1 - u \pmod{2, \omega}\]
ist, und $S(u)$ verschwindet nur, wenn
\[w \equiv 0, 1 \pmod{2, \omega}\]
ist.

Da jede gebrochene Zahl in $\Omega$ durch Multiplikation mit einer ganzen rationalen Zahl in eine ganze Zahl verwandelt werden kann, so ergibt sich aus \S 156, (6) der Satz:

\begin{satz}
Ist $\eta$ irgend eine gebrochene Zahl des K\"orpers $\Omega$, so ist $S(\eta)$ eine algebraische Zahl, und wenn insbesondere $\eta$ so dargestellt werden kann, da\ss{} sein Nenner relativ prim zu 2 ist, so ist $S(\eta)$ eine ganze algebraische Zahl.
\end{satz}

Es sei jetzt
\begin{equation}\mu = y + x\omega\end{equation}
eine ganze Zahl in $\Omega$, in der $x, y$ ganze rationale Zahlen sind, und
\[\bar{\mu} = y + x\bar{\omega}\]
die zu $\mu$ konjugierte Zahl. Ferner
\begin{equation}m = \mu\bar{\mu} = y^2 - bxy + acx^2\end{equation}
die Norm von $\mu$, die wir als ungerade voraussetzen.

\section*{\S 158. Zerlegung der Funktion $A(x)$.}
\addcontentsline{toc}{section}{\S 158. Zerlegung der Funktion $A(x)$.}

Es seien jetzt
\begin{equation}\mu = y + x\omega, \quad \mu_1 = y_1 + x_1\omega\end{equation}
zwei ungerade ganze Zahlen in $\Omega$ von der Form \S 157, (6).

Es sei $m$ der gr\"o\ss{}te gemeinschaftliche Idealteiler von $\mu$ und $\mu_1$.

Setzt man also
\[\mu = m\alpha, \quad \mu_1 = m\alpha_1,\]
so sind $\alpha$ und $\alpha_1$ zwei \"aquivalente Ideale ohne gemeinsamen Teiler.

Die den beiden Zahlen $\mu, \mu_1$ entsprechenden Funktionen $A(x)$ bezeichnen wir mit
\[A_\mu(x), \quad A_{\mu_1}(x);\]
und fragen, welche gemeinschaftliche Wurzeln diese Funktionen haben, wann also die Gleichung:
\begin{equation}S\left(\frac{\varrho}{\mu}\right) = S\left(\frac{\varrho_1}{\mu_1}\right)\end{equation}
erf\"ullt sein kann, wenn $\varrho$ nach dem Modul $\mu$, $\varrho_1$ nach dem Modul $\mu_1$ genommen ist.

Da die Kongruenz
\[\frac{2\varrho}{\mu} \equiv \frac{2\varrho_1}{\mu_1} \pmod{2, \omega}\]
nicht m\"oglich ist, was man ganz wie oben (S.~589) zeigt, so ist f\"ur die Gleichung (3) notwendig und hinreichend:
\begin{equation}\varrho\mu_1 \equiv 0 \pmod{2, \omega}.\end{equation}

Daraus folgt:
\begin{equation}\varrho_1 \equiv \mu\xi \pmod{\mu_1},\end{equation}
und dies ist nur dann m\"oglich, wenn
\begin{equation}\varrho \equiv 0 \pmod{\alpha}, \quad \varrho_1 \equiv 0 \pmod{\alpha_1}.\end{equation}

Sind umgekehrt die Bedingungen (5) und (6) erf\"ullt, so ist:
\[\frac{\varrho}{\mu} = \frac{\varrho_1}{\mu_1}\]
eine ganze Zahl und die Gleichung (3) ist befriedigt.

Man nehme nun eine durch $\alpha$ teilbare ganze Zahl in $\Omega$:
\begin{equation}\beta = \alpha\gamma,\end{equation}
worin $\gamma$ relativ prim zu $\mu$ und $\mu_1$ ist.

Wenn dann $\varrho$ durch $\alpha$ teilbar ist, so kann man eine ganze Zahl $\xi$ nach dem Modul $m$ aus der Kongruenz
\[\varrho = \mu\xi \pmod{\mu}\]
bestimmen (Bd.~II, \S 166, 7.); dann ist
\[\frac{\varrho}{\mu} - \xi = \frac{\beta\xi}{\mu}\]
eine durch $\alpha\gamma$ und durch kein anderes Ideal teilbare ganze Zahl, und aus (5) ergibt sich
\[\varrho_1 = \mu_1\xi \pmod{\mu_1}.\]

\begin{satz}
Wir erhalten also die gemeinschaftlichen Wurzeln von $A_\mu(x)$ und $A_{\mu_1}(x)$ in der Form:
\begin{equation}S\left(\frac{\beta\xi}{\mu}\right),\end{equation}
worin $\xi$ ein vollst\"andiges Restsystem nach dem Modul $m$ durchl\"auft.
\end{satz}

Hierin kann $m$ jedes beliebige ungerade Ideal in $\Omega$ bedeuten.

Denn man kann zu jedem solchen Ideal zwei Zahlen $\mu, \mu_1$ von der Form (2) w\"ahlen und dann $\beta$ unter den \"aquivalenten Zahlen so, da\ss{} $\alpha$ ungerade und relativ prim zu $\omega$ wird.

\section*{\S 159. Primideale.}
\addcontentsline{toc}{section}{\S 159. Primideale.}

\begin{satz}
Ist $\eta$ eine gebrochene Zahl des K\"orpers $\Omega$, die in reduzierter Form den ungeraden Idealnenner $\mathfrak{a}$ hat, und ist $m$ ein zu $\mathfrak{a}$ relativ primes Ideal, so ist die ganze Zahl $S(\eta)$ relativ prim zu $m$.
\end{satz}

Denn man nehme in $\Omega$ eine durch $\mathfrak{a}$ teilbare, zu $2m$ teilerfremde ganze Zahl $\mu$ an. Dann ist $\mu\eta$ eine ganze Zahl.

In dem Produkt \S 157, (23)
\begin{equation}\prod_{\varrho} S\left(\eta + \frac{\varrho}{\mu}\right)\end{equation}
kommt eine Zahl $2\varrho$ vor, die mit $\mu\eta$ nach dem Modul $\mu$ kongruent ist. Folglich ist $S(\eta)$ ein Teiler von $\mu$ und mithin relativ prim zu $m$, wie bewiesen werden sollte.

Nehmen wir jetzt an, da\ss{} in der Formel (11), \S 158
\begin{equation}r_{\mathfrak{p}} = \mu \prod_{\xi} S\left(\frac{\beta\xi}{\mu}\right)\end{equation}
$m$ ein Primideal des K\"orpers $\Omega$ sei und setzen demgem\"a\ss{} $m = \mathfrak{p}$, so da\ss{} $\xi$ ein volles Restsystem nach dem Modul $\mathfrak{p}$ mit Ausschlu\ss{} der Null durchl\"auft.

Nehmen wir dann $\mu$ so an, da\ss{} es nur durch die erste Potenz von $\mathfrak{p}$ teilbar und durch ein beliebig gew\"ahltes anderes Primideal $\mathfrak{p}'$ nicht teilbar ist, so kann $r_{\mathfrak{p}}$ nach dem Satz 6. nicht durch $\mathfrak{p}'$ teilbar sein.

Nun ist $r_{\mathfrak{p}}$ nur von dem Ideal $\mathfrak{p}$ abh\"angig und unabh\"angig davon, wie im \"ubrigen die Zahl $\mu$ genommen ist.

Folglich ist $r_{\mathfrak{p}}$ nach dem Satz 6. durch kein von $\mathfrak{p}$ verschiedenes Primideal teilbar.

Es kann aber auch $\mathfrak{p}$ nicht in einer h\"oheren als der ersten Potenz in $r_{\mathfrak{p}}$ aufgehen, denn die Faktoren des Produktes (1) kommen alle unter den Faktoren des Produktes (2) vor, und folglich ist $\mu$ durch $r_{\mathfrak{p}}$ teilbar.

Demnach k\"onnen wir geradezu
\begin{equation}\mu = \mathfrak{p}\end{equation}
setzen, und da $r_{\mathfrak{p}}$ eine Zahl im K\"orper $\Omega$ ist, so folgt:

\begin{satz}
Jedes ungerade Primideal des K\"orpers $\Omega$ ist in dem K\"orper $\mathfrak{K}$ ein Hauptideal.
\end{satz}

\section*{\S 160. Primideale ersten Grades in $\mathfrak{T}_u$.}
\addcontentsline{toc}{section}{\S 160. Primideale ersten Grades in $\mathfrak{T}_u$.}

Es kommt nun vor allem darauf an, die Primzahlen $p$ aufzusuchen, die im K\"orper $\Omega$ in Primideale ersten Grades zerlegbar sind.

Nach \S 122 und \S 156 sind das die Primzahlen $p$, die durch die Hauptform der Diskriminante $D = 4\Delta$ oder $16\Delta$ darstellbar sind.

Diese sind von der Form:
\begin{enumerate}
\item $p = y^2 - \Delta x^2$, $\Delta \equiv 0 \pmod{4}$ oder $\equiv 1 \pmod{8}$,
\item $p = y^2 - 4\Delta x^2$, $4\Delta \equiv 5 \pmod{8}$.
\end{enumerate}

Ist $4\Delta \equiv 1 \pmod{8}$, so mu\ss{} der dritte Koeffizient der Form $(a, 2b, c)$ gerade sein.

Ist $\Delta \equiv 0 \pmod{4}$, so k\"onnen wir $c$ gerade annehmen, indem wir n\"otigenfalls zu der Parallelform $(a, b + 2a, a + b + c)$ \"ubergehen.

Zerf\"allt also $p$ in die beiden Primfaktoren $\pi, \pi'$, so ist
\begin{equation}\pi = y + x\sqrt{\Delta}, \quad \pi' = y - x\sqrt{\Delta},\end{equation}
und um die Multiplikationsformel \S 157, (10), (17) auf $\mu = \pi$ anzuwenden, haben wir im letzteren Falle $\xi$ durch $2\pi$ zu ersetzen.

Wir k\"onnen also in allen F\"allen $cx$ gerade und daher $\varepsilon = \pm 1$ annehmen.

\section*{\S 161. Zahlgruppen und Idealgruppen.}
\addcontentsline{toc}{section}{\S 161. Zahlgruppen und Idealgruppen.}

Wie in \S 98 und \S 106 bezeichne ich mit $\mathfrak{O}$ die Gesamtheit der ganzen und gebrochenen Zahlen des quadratischen K\"orpers $\Omega$, nach Ausschlu\ss{} aller Zahlen, die zu einem beliebig angenommenen $\mathfrak{E}$ nicht teilerfremd sind.

$\mathfrak{E}$ kann eine rationale Zahl, eine Zahl in $\Omega$ oder auch ein Ideal sein. Ich will es den Exkludenten nennen.

In den drei Abhandlungen "uber Zahlgruppen in algebraischen K\"orpern\footnote{Mathematische Annalen Bd. 48, 49, 50.} habe ich gewisse Systeme von Zahlen aus $\mathfrak{O}$ unter dem Namen Zahlgruppen zusammengefa\ss{}t, die dadurch definiert waren, da\ss{} das Produkt und der Quotient je zweier Zahlen einer solchen Gruppe in derselben Gruppe enthalten waren.

Eine solche Zahlgruppe kann niemals die Zahl Null enthalten, enth\"alt aber sicher die Zahl $1$.

Im gleichen Sinne wie von Zahlgruppen k\"onnen wir auch von Idealgruppen reden, die ebenfalls ihren Exkludenten haben k\"onnen.

Es kommen hierbei nat\"urlich auch gebrochene Ideale vor, und man bedient sich dabei nicht ohne Nutzen der Darstellungsweise durch die Funktionale (Bd.~I, \S 169).

Die Gesamtheit der ganzen und gebrochenen Zahlen des K\"orpers $\Omega$, der wir einen beliebigen Exkludenten $\mathfrak{E}$ beilegen, bildet die Zahlgruppe $\mathfrak{O}$ und die Gesamtheit der Ideale $\mathfrak{O}$ in dem gleichen Sinne eine Idealgruppe.

Ebenso bildet die Gesamtheit der numerischen Einheiten eine Zahlgruppe $\mathfrak{E}$ und die Gesamtheit der funktionalen Einheiten eine Idealgruppe $\mathfrak{E}$.

Die Multiplikation zweier Gruppen $\mathfrak{A}$ und $\mathfrak{B}$ zu dem Produkt $\mathfrak{A} \cdot \mathfrak{B}$ geschieht dadurch, da\ss{} man jedes Element von $\mathfrak{A}$ mit jedem Element von $\mathfrak{B}$ multipliziert.

Ist $\mathfrak{A}$ eine Zahlgruppe, so ist $\mathfrak{E}\mathfrak{A}$ eine Idealgruppe. Diese enth\"alt aber nur Hauptideale und kann daher Hauptidealgruppe genannt werden.

Wenn die ganze Gruppe $\mathfrak{O}$ in eine endliche Anzahl von Nebengruppen nach $\mathfrak{A}$ zerf\"allt:
\begin{equation}\mathfrak{O} = \mathfrak{A}\omega_1 + \mathfrak{A}\omega_2 + \cdots + \mathfrak{A}\omega_h,\end{equation}
so hei\ss{}t die Zahl $h$ der Index von $\mathfrak{A}$ und wird mit
\begin{equation}h = (\mathfrak{O}, \mathfrak{A})\end{equation}
bezeichnet (wie in \S 100).

\section*{\S 162. Die durch ein Ideal teilbaren Ideale der Hauptklasse.}
\addcontentsline{toc}{section}{\S 162. Die durch ein Ideal teilbaren Ideale der Hauptklasse.}

Es soll jetzt unter $\Omega$ der quadratische K\"orper mit negativer Grundzahl $\Delta$ verstanden werden.

In diesem K\"orper sei $\mathfrak{A}$ eine Kongruenzgruppe mit dem Modul $M$ und dem Exkludenten $\mathfrak{E}$.

Die Gruppe $\mathfrak{A}$ enth\"alt, wie wir gesehen haben, unendlich viele ganze Zahlen.

Wir fragen nach der Anzahl $\tau(t)$ der ganzen Zahlen in $\mathfrak{A}$, die durch irgend ein gegebenes ganzes zu $M$ teilerfremdes Ideal $\mathfrak{m}$ teilbar sind, deren Norm nicht gr\"o\ss{}er als die positive Zahl $t$ ist.

Es sei $(\omega_1, \omega_2)$ eine Basis von $\mathfrak{m}$, also nach \S 91, (11):
\begin{equation}\mathfrak{m} = (a, \frac{b + \sqrt{\Delta}}{2}),\end{equation}
\[b \equiv \sqrt{\Delta} \pmod{2a}, \quad \Delta = b^2 - 4ac, \quad N(\mathfrak{m}) = ac.\]

Da $\mathfrak{m}$ relativ prim zu $M$ vorausgesetzt ist, so k\"onnen wir nach \S 161, (6) in jeder in $\mathfrak{A}$ vorkommenden Zahlklasse nach $M$ eine durch $\mathfrak{m}$ teilbare Zahl $\alpha_0$ bestimmen.

\section*{\S 163. Die Dirichletschen Summen.}
\addcontentsline{toc}{section}{\S 163. Die Dirichletschen Summen.}

Wir verstehen jetzt unter $\mathfrak{O}$ die Idealgruppe der s\"amtlichen Ideale in $\Omega$ (mit R\"ucksicht auf den Exkludenten $\mathfrak{E}$); $\mathfrak{A}$ sei eine Kongruenzgruppe und
\begin{equation}\mathfrak{A} = \mathfrak{E}\mathfrak{A}\end{equation}
die Gruppe der entsprechenden Hauptideale.

Wir zerlegen $\mathfrak{O}$ nach \S 161, (3) in die Klassen nach $\mathfrak{A}$:
\begin{equation}\mathfrak{O} = \mathfrak{A}_1 + \mathfrak{A}_2 + \cdots + \mathfrak{A}_h,\end{equation}
und betrachten die Summen:
\begin{equation}S_i = \sum_{\mathfrak{a}_i} \frac{1}{N(\mathfrak{a}_i)^s},\end{equation}
worin $s$ eine positive Variable $> 1$ bedeutet, und $\mathfrak{a}_i$ die s\"amtlichen ganzen Ideale einer Klasse $\mathfrak{A}_i$ durchl\"auft.

\section*{\S 164. Der Klassenk\"orper.}
\addcontentsline{toc}{section}{\S 164. Der Klassenk\"orper.}

Zu der Zahlgruppe $\mathfrak{A}$ in dem quadratischen K\"orper $\Omega$ soll ein algebraischer K\"orper \"uber $\Omega$ existieren, den man mit $\Omega(\mathfrak{A})$ oder $\mathfrak{K}(\mathfrak{A})$ bezeichne, von dem wir nur voraussetzen wollen:

\begin{satz}[Definition des Klassenk\"orpers]
Die Primideale $\mathfrak{p}_1$ ersten Grades der Hauptklasse $\mathfrak{A}_0$, und nur diese, sollen im K\"orper $\Omega(\mathfrak{A})$ wieder in Primideale ersten Grades zerfallen.
\end{satz}

Es sind also nur die Primideale in $\Omega$, die wir vorhin mit $\mathfrak{p}_1$ bezeichnet haben, die in $\mathfrak{K}$ in Primideale ersten Grades zerfallen.

Eine beliebige endliche Anzahl von Primzahlen k\"onnen dabei ausgenommen sein. Sie werden dann in den Exkludenten genommen.

Ist $n$ der relative Grad des K\"orpers $\mathfrak{K}$ \"uber $\Omega$ und $p$ eine nat\"urliche Primzahl, in der die Ideale $\mathfrak{p}_1$ und $\mathfrak{P}$ in $\Omega$ und $\mathfrak{K}$ aufgehen, und bezeichnet man mit $N_\Omega$, $N_\mathfrak{K}$ die absoluten Normen im K\"orper $\Omega$ und $\mathfrak{K}$, so ist
\[N\mathfrak{P} = N_\mathfrak{K}\mathfrak{P} = p^{f_1} = p^n\]
und folglich mu\ss{} $\mathfrak{p}_1$ im K\"orper $\mathfrak{K}$ in $n$ Primideale $\mathfrak{P}$ zerfallen:
\begin{equation}\mathfrak{p}_1 = \mathfrak{P}_1 \mathfrak{P}_2 \cdots \mathfrak{P}_n.\end{equation}

Die Primideale, die in den Grundzahlen der K\"orper $\Omega$ oder $\mathfrak{K}$ aufgehen, sollen immer ausgeschlossen sein.

Dann sind die $\mathfrak{P}_1, \mathfrak{P}_2, \ldots, \mathfrak{P}_n$ voneinander verschieden.

Ob ein solcher K\"orper $\mathfrak{K}$ existiert, bleibt vorl\"aufig unentschieden.

Wir wollen untersuchen, was aus der Definition geschlossen werden kann.

Da in jedem Ideal $\mathfrak{p}_1$ $n$ Ideale $\mathfrak{P}$ aufgehen, so ist
\begin{equation}n \sum \frac{1}{N(\mathfrak{p}_1)^s} = \sum \frac{1}{N(\mathfrak{P})^s},\end{equation}
worin sich das erste Produkt auf alle Ideale $\mathfrak{p}_1$ des K\"orpers $\Omega$, das zweite auf alle Ideale $\mathfrak{P}$ ersten Grades des K\"orpers $\mathfrak{K}$ erstreckt,

und nach dem f\"ur jeden beliebigen K\"orper g\"ultigen Satz I, \S 197 des II.~Bandes hat
\begin{equation}\lim_{s \to 1} (s-1) \sum \frac{1}{N(\mathfrak{P})^s}\end{equation}
f\"ur $s = 1$ einen von Null verschiedenen endlichen Grenzwert, der dort auf die Klassenzahl im K\"orper $\mathfrak{K}$ zur\"uckgef\"uhrt ist.

Setzen wir zur Abk\"urzung
\[P = \prod_{\mathfrak{p}_1} \frac{1}{1 - N(\mathfrak{p}_1)^{-s}},\]
so ist nach (2) und (3)
\begin{equation}\lim_{s \to 1} (s-1)^n P^n \text{ endlich und von Null verschieden,}\end{equation}
und nach \S 163, 6.:
\begin{equation}\lim_{s \to 1} (s-1) P \text{ endlich,}\end{equation}
und wenn man hieraus $P$ eliminiert, so folgt:
\[\lim_{s \to 1} (s-1)^{n-1} \text{ endlich.}\]

Also ist
\begin{equation}n \geq h.\end{equation}

Daraus folgt der Satz:

\begin{satz}
Der Relativgrad $n$ des Klassenk\"orpers kann nicht kleiner sein als die Klassenzahl $h$ der Gruppe.
\end{satz}

\section*{\S 165. Primideale in den Klassen.}
\addcontentsline{toc}{section}{\S 165. Primideale in den Klassen.}

Wir kehren jetzt zu den Funktionen $Q$ zur\"uck, die wir in \S 163, (16) durch unendliche Reihen und in \S 163, (18) durch unendliche Produkte:
\[Q_\chi = \prod_{\mathfrak{p}} \frac{1}{1 - \chi(\mathfrak{p}) N(\mathfrak{p})^{-s}}\]
dargestellt haben, worin $\chi$ einen der Charaktere der Gruppe bedeutet und $\mathfrak{p}$ alle Primideale in $\mathfrak{O}$ durchl\"auft.

Den Logarithmus dieses Produktes entwickeln wir in folgender Weise:
\begin{equation}\log Q_\chi = \sum_{\mathfrak{p}} \frac{\chi(\mathfrak{p})}{N(\mathfrak{p})^s} + \frac{1}{2} \sum_{\mathfrak{p}} \frac{\chi(\mathfrak{p}^2)}{N(\mathfrak{p})^{2s}} + \frac{1}{3} \sum_{\mathfrak{p}} \frac{\chi(\mathfrak{p}^3)}{N(\mathfrak{p})^{3s}} + \cdots\end{equation}

Das Vielfache von $2\pi i$, das in dem $\log Q$ nicht n\"aher definiert ist, brauchen wir nicht zu kennen, da es sich mit $s$ nur unstetig \"andern k\"onnte, und da rechts und links stetige Funktionen von $s$ stehen, solange $s > 1$ ist, so \"andert sich dieses Vielfache \"uberhaupt nicht mit $s$.

Es sei $\mathfrak{A}_i$ eine der Klassen der Idealgruppe $\mathfrak{A}$ und $\mathfrak{A}_i'$ die entgegengesetzte Klasse.

Dann ist nach (12), (13), (14), \S 163:
\[\sum_\chi \chi(\mathfrak{A}_i') \log Q_\chi = h \sum_{\mathfrak{p} \in \mathfrak{A}_i} \frac{1}{N(\mathfrak{p})^s} + \frac{h}{2} \sum_{\mathfrak{p}^2 \in \mathfrak{A}_i} \frac{1}{N(\mathfrak{p})^{2s}} + \frac{h}{3} \sum_{\mathfrak{p}^3 \in \mathfrak{A}_i} \frac{1}{N(\mathfrak{p})^{3s}} + \cdots\]

\section*{\S 166. Primideale in den Idealklassen.}
\addcontentsline{toc}{section}{\S 166. Primideale in den Idealklassen.}

Nehmen wir als Gruppe zun\"achst die Gesamtheit der Zahlen $\mathfrak{O}$ des K\"orpers $\Omega$ ohne die Null, so sind die Klassen die Idealklassen des K\"orpers $\Omega$, die, wie wir gesehen haben, den Klassen quadratischer Formen der Diskriminante $\Delta$ eindeutig zugeordnet werden k\"onnen.

Nehmen wir als Gruppe die Ordnung $\mathfrak{O}' = [\mathfrak{O}]$ mit dem F\"uhrer $\mathfrak{f}$, die durch die Kongruenzbedingung definiert ist, da\ss{} jede Zahl in $\mathfrak{O}'$ nach dem Modul $\mathfrak{f}$ mit einer rationalen Zahl kongruent sein soll, so haben wir eine Kongruenzgruppe, deren Exkludent und Modul $= \mathfrak{f}$ ist.

Die Klassen entsprechen den Formenklassen der Diskriminante $D = \mathfrak{f}^2 \Delta$.

Der zweite Fall schlie\ss{}t den ersten in sich (f\"ur $\mathfrak{f} = 1$).

Der Klassenk\"orper ist hier der Klassenk\"orper $\Omega(D)$, den wir zum Unterschied von allgemeineren als Ordnungsk\"orper bezeichnen wollen.

Denn ist $(k)$ eine Klasseninvariante, so ist
\[(k)^p = (k) \pmod{\mathfrak{P}}\]
nur dann befriedigt, wenn $p$ durch die Hauptklasse der Diskriminante $D$ darstellbar ist.

Dann zerf\"allt $p$ in $\Omega$ in zwei konjugierte Primideale $\mathfrak{p}, \mathfrak{p}'$, die durch Zahlen in $\mathfrak{O}$ oder in $\mathfrak{O}'$ darstellbar, also Hauptideale sind.

\section*{\S 167. Primzahlen in Linearformen.}
\addcontentsline{toc}{section}{\S 167. Primzahlen in Linearformen.}

Wir definieren nun eine Zahlgruppe $\mathfrak{A}$ in $\Omega$ folgenderma\ss{}en:

Es sei $m$ ein ganzes Ideal in $\Omega$, und $\mathfrak{A}$ bestehe aus allen ganzen und gebrochenen Zahlen $\xi$ in $\Omega$, die zu $m$ teilerfremd sind und der Bedingung gen\"ugen:
\begin{equation}\xi \equiv 1 \pmod{m}.\end{equation}

Nehmen wir die Primteiler von $m$ in den Exkludenten, und bezeichnen mit $\mathfrak{O}$ die Gruppe der Zahlen in $\Omega$, so zerf\"allt $\mathfrak{O}$ nach dem Modul $m$ in $\psi(m)$ Zahlklassen, worin $\psi(m)$ die Bedeutung wie in Bd.~II, \S 168 hat, n\"amlich, wenn $\mathfrak{p}$ die Primteiler von $m$ durchl\"auft:
\begin{equation}\psi(m) = N(m) \prod_{\mathfrak{p}|m} \left(1 - \frac{1}{N(\mathfrak{p})}\right).\end{equation}

In die Gruppe $\mathfrak{A}$ geh\"oren nicht blo\ss{} einzelne Zahlen, sondern Zahlklassen nach dem Modul $m$.

\section*{\S 168. Reduktion der Klassengleichung in den Kreisteilungsk\"orpern.}
\addcontentsline{toc}{section}{\S 168. Reduktion der Klassengleichung in den Kreisteilungsk\"orpern.}

Wir definieren eine weitere Gruppe $\mathfrak{A}$ durch folgende Bestimmung:

Es seien $\mathfrak{f}$, $m$ zwei ganze rationale Zahlen, deren Primfaktoren wir in den Exkludenten $\mathfrak{E}$ aufnehmen.

$\mathfrak{O}$ sei wie in \S 98 die Gruppe der Zahlen in $\Omega$, und $\mathfrak{O}'$ die Gruppe der Zahlen der Ordnung $[\mathfrak{f}]$, und $\mathfrak{A}$ bestehe aus allen ganzen und gebrochenen Zahlen $\xi$ in $\mathfrak{O}'$, die der Kongruenzbedingung
\begin{equation}N(\xi) \equiv 1 \pmod{m}\end{equation}
gen\"ugen.

Zun\"achst haben wir die Klassenzahl
\begin{equation}h' = (\mathfrak{O}', \mathfrak{E}\mathfrak{A})\end{equation}
zu bestimmen.

Der erste Faktor ist die Klassenzahl der quadratischen Formen der Ordnung $\mathfrak{O}'$ [\S 100, (3)], die wir mit $H'$ bezeichnen.

Es ist also
\begin{equation}h' = \frac{H' \cdot (\mathfrak{O}', \mathfrak{A})}{(\mathfrak{E}\mathfrak{A}, \mathfrak{A})}\end{equation}
(nach \S 100, 13.).

Es bleibt noch $(\mathfrak{O}', \mathfrak{A})$ zu bestimmen.

\end{document}
